\section{Full multiplicity from transverse jets along a binary plane}

\subsection{O-020: arbitrary second rows with third row at most two}
Let $r\geq3$ and $n\geq2r-4$. For every outer degree $d$
and every partition $\lambda\vdash nd$ with
$\ell(\lambda)\leq r$ and $\lambda_3\leq2$,
\[
                         K_n(\lambda)=a_\lambda(d[n]).
                                                               \tag{O-020.1}
\]
Missing third parts are zero. There is no bound on
$\lambda_2$ and no upper bound on $d$.
The geometric input is dominance of the determinant projection
onto the entire coefficient space of degree-$n$ forms whose
total degree outside a fixed binary plane is at most two.
Higher transverse coefficients remain unrestricted.

Write the variables as $x,y,z_1,\ldots,z_R$, where $R=r-2$.
The target coefficient space is
\[
 \mathcal T_{n,2}=
 \operatorname{Sym}^n\mathbb C^2
 \oplus\left(\operatorname{Sym}^{n-1}\mathbb C^2
                               \otimes\mathbb C^R\right)
 \oplus\left(\operatorname{Sym}^{n-2}\mathbb C^2
                               \otimes\operatorname{Sym}^2\mathbb C^R\right).
                                                               \tag{O-020.2}
\]
Its dimension is $(n+1)+Rn+\binom{R+1}{2}(n-1)$.
The map from all $n$-by-$n$ linear pencils in these variables
to $\mathcal T_{n,2}$, obtained by deleting terms of
total $z$-degree at least three from the determinant,
has a surjective differential at an explicit integral pencil.

\paragraph{Edge-disjoint spanning trees.}
We first fix $R$ edge-disjoint spanning trees on $n$ vertices,
possible because $n\geq2R$. The following elementary cyclic
construction specifies the needed edge colors.
On vertices $\mathbb Z/(2R)$, for $0\leq j<R$ take
the spanning path
\[
 j,\quad (j-i,j+i)_{1\leq i\leq R-1},\quad j-R.
\]
Call this path $T_{j+1}$. The paired middle list is empty
when $R=1$. These paths are edge-disjoint. Indeed the edges
in the $j$-th path have endpoint sums $2j-1$ or $2j$
modulo $2R$. Each odd sum gives a perfect matching;
each even sum gives a matching excluding its two loop
solutions. As $j$ varies these sums exhaust their respective
parities without repetition. This proves edge-disjointness;
the paths in fact partition all edges of the complete graph
on $2R$ vertices.
For each additional vertex, attach it in tree $T_a$ to a
different one of the original vertices for each $a$.
Those added edges are distinct across colors and make the
new vertices leaves, so every $T_a$ remains a spanning tree.

Relabel all vertices $1,\ldots,n$ and orient every edge
from its smaller endpoint to its larger endpoint. For
$i<j$, put $u_{ij}=z_a$ when $\{i,j\}$ belongs to
tree $T_a$, and put $u_{ij}=0$ otherwise. The edge-disjointness
makes this definition unambiguous. Let $U=(u_{ij})$ be
strictly upper triangular, choose distinct integers
$a_1,\ldots,a_n$, and define
\[
 \ell_i=y-a_ix,\qquad \Delta=\prod_{i=1}^n\ell_i,
             A_0=\operatorname{diag}(\ell_1,\ldots,\ell_n)-U.
                                                               \tag{O-020.3}
\]
Then $\det A_0=\Delta$.

\paragraph{Two binary spanning facts.}
The $n$ polynomials $\Delta/\ell_i$ form a basis of the
binary forms of degree $n-1$. After setting $x=1$, their
evaluation matrix at $y=a_j$ is diagonal with nonzero
diagonal entries.
For the edges of any spanning tree $T$ on these vertices,
the $n-1$ polynomials
\[
                           \frac{\Delta}{\ell_i\ell_j},
                                   \qquad\{i,j\}\in T,          \tag{O-020.4}
\]
form a basis of binary forms of degree $n-2$.
To see this explicitly, send a polynomial $p(y)$ of degree
at most $n-2$ to the vector
\[
                \left(\frac{p(a_i)}{\Delta'(a_i)}\right)_{i=1}^n.
\]
Lagrange interpolation shows that its coordinates sum to
zero; the map is injective and both spaces have dimension
$n-1$, so it identifies this polynomial space with the
zero-sum hyperplane. The image of
$\Delta/((y-a_i)(y-a_j))$ is
$(e_i-e_j)/(a_i-a_j)$. The edge differences of a tree
form a basis of the zero-sum hyperplane, proving the claim.
Rehomogenization gives (O-020.4).

\begin{proof}[Proof of the surjective differential]
The cofactor derivative formula is
\[
             d(\det)_{A_0}(H E_{ij})
                         =H\,\Delta\,(A_0^{-1})_{ji}.
                                                               \tag{O-020.5}
\]
The expression is polynomial although it is convenient to
compute the inverse over $\mathbb C(x,y,z)$.
Since $U$ is strictly upper triangular, an inverse entry
is a finite sum over increasing paths. In particular,
\[
 (A_0^{-1})_{ii}=\ell_i^{-1},\qquad
 (A_0^{-1})_{ij}
       =\frac{u_{ij}}{\ell_i\ell_j}
                    +\text{terms of total $z$-degree at least two}
                       \quad(i<j).                    \tag{O-020.6}
\]
Every extra path edge contributes another $z$ factor.
Multiplication by $\Delta$ clears the denominators,
because a strictly increasing path visits each vertex once.

Vary diagonal entries first. Taking $H=x$ or $H=y$
in (O-020.5) spans every binary degree-$n$ form:
the polynomials $\Delta/\ell_i$ span degree $n-1$,
and multiplication by $x,y$ spans the next degree.
Taking $H=z_a$ on diagonal entries spans independently
every term linear in the $z$ variables with an arbitrary
binary coefficient of degree $n-1$.

Now fix a color $a$ and an edge $\{i,j\}$ of its tree,
with $i<j$. Vary the opposite, lower-triangular entry
by $H=z_b$, for any $1\leq b\leq R$.
Equations (O-020.5)--(O-020.6), after deleting transverse
degree at least three, give exactly
\[
                          z_az_b\frac{\Delta}{\ell_i\ell_j}.
                                                               \tag{O-020.7}
\]
The terms from longer paths have degree at least three
after multiplication by $z_b$ and disappear from the
target projection. As the edges range over tree $T_a$,
(O-020.4) says that the binary coefficient in (O-020.7)
can be arbitrary. As $a,b$ vary, their products span
all quadratic forms in the transverse variables.

These directions therefore span every summand of
(O-020.2). Diagonal directions control transverse orders
zero and one, while the lower-triangular directions just
used have no contributions in those orders. Edge colors
are unambiguous, and all matrix-entry coefficients are
independent parameters. Thus the differential is surjective,
with no requirement to fix higher transverse terms.
The algebraic lowest-homogeneous-term argument in M-013
now proves dominance of this coefficient projection.
\end{proof}

\begin{proof}[Proof of the multiplicity equality]
Let $w$ be any degree-$d$ highest-weight vector of weight
$\lambda$ in the determinant ideal. It depends only on
the first $r$ coordinates. By the $j=3$ support bound
of F-016, every coefficient generator used by $w$ has
total degree at most $\lambda_3\leq2$ in variables
$3,\ldots,r$. Hence $w$ is a polynomial in the
coordinates of $\mathcal T_{n,2}$.
Every linear pencil determinant belongs to $\Omega_n$,
so this polynomial vanishes after the dominant projection
just proved. Its pullback is injective, so $w=0$.
This holds for every highest-weight copy, and therefore
the ideal multiplicity space vanishes entirely. The
ambient quotient then gives (O-020.1).
\end{proof}

For permanent shapes $r\leq m^2$, the condition
$n\geq2m^2-4$ is enough for this full equality whenever
$\lambda_3\leq2$, with the one- and two-row cases already
covered by O-006. This conclusion is stronger than a
comparison with a particular permanent multiplicity: every
ambient copy survives, even with an arbitrarily large second
row. It is a different sufficient condition from M-013's
bounded second-row jet theorem.

\paragraph{Geometric dictionary.}
The target is the degree-$n$ part of the quotient by
$(z_1,\ldots,z_R)^3$, the second infinitesimal neighborhood
transverse to the binary coordinate plane. It is a finite
locally free coefficient module. The colored graph, diagonal
linear forms, and cofactor derivatives define integral
morphisms. Distinctness of the $a_i$ and the nonzero
Vandermonde factors are used after base change to
$\mathbb Q$ or $\mathbb C$; no equal-rank assertion in
every residue characteristic is inferred.
The suffix support is a torus/comodule calculation under
SGA~3 I, Propositions~4.7.2--4.7.3, pp.~34--36,
and the Cauchy and Schur functors of Macdonald I(4.3),
p.~63, and Appendix~A(5.3)--(5.6), pp.~155--156.
Dominance injects the actual coordinate algebra of the
detected coefficients, which makes the whole highest-weight
ideal component vanish. Neither exact lifting of each
truncated form nor containment of that entire linear space
in the determinant closure is asserted.
